\section{Applications of the discrete Fourier transform}

\begin{correction}{exo-transforme-sup-comp}
Compact support and Cauchy--Schwarz give $\wh f\in L^1$. Fourier inversion
therefore defines a continuous representative. Differentiating under the
integral, with the bound $|\xi|^n\leq A^n$, gives
\[
 f^{(n)}(t_0)=\frac1{2\pi}\int_{-A}^A
             (\imath\xi)^n\wh f(\xi)e^{\imath\xi t_0}\d\xi.
\]
For every $z\in\CC$, the exponential series is dominated on this interval by
$e^{A|z-t_0|}|\wh f(\xi)|$, an integrable function. Thus
\[
 \frac1{2\pi}\int_{-A}^A\wh f(\xi)e^{\imath\xi z}\d\xi
 =\sum_{n=0}^\infty\frac{f^{(n)}(t_0)}{n!}(z-t_0)^n.
\]
This constructs an entire extension of $f$. If its real restriction vanishes
on an open interval, all derivatives at an interior point vanish, so the
series shows it is identically zero. If $f$ also has compact support,
its continuous representative vanishes outside that support and hence on an
open interval. Therefore $f=0$, proving the proposition.
\end{correction}

\begin{correction}{exo-resol-eq-heat-diff-finies}
\begin{enumerate}
\item Use the forward difference $(u(t+h,x)-u(t,x))/h$ in time and the
centered difference $(u(t,x+d)-2u(t,x)+u(t,x-d))/d^2$ in space.
The local consistency errors are $O(h)$ and $O(d^2)$ for smooth $u$.
\item The kernel is $g[0]=1-2s$, $g[\pm1]=s$. One direct step costs $O(N)$;
an FFT at every step would cost $O(N\log N)$. If only step $m$ is wanted,
two FFTs and the multipliers $\wh g[k]^m$ avoid computing all intermediate
steps. Program~\ref{listing-resolexplicite} implements the local stencil.
\begin{listing}[H]
\PythonCode{listing-resolexplicite}
\caption{Explicit heat evolution by a three-point stencil.}
\label{listing-resolexplicite}
\end{listing}
\item The multiplier is
\[
 r_k=1-4s\sin^2(\pi k/N).
\]
The DFT diagonalizes every power of the iteration. Boundedness for every
initial vector is equivalent to $|r_k|\leq1$ at every frequency, hence
$2s\mu_N\leq1$, where
\[
 \mu_N=\max_k\sin^2(\pi k/N)=
 \begin{cases}1,&N\text{ even},\\\cos^2(\pi/(2N)),&N\text{ odd}.\end{cases}
\]
This gives the stated bounds. The name CFL abbreviates
Courant--Friedrichs--Lewy. For $s\leq1/2$, the stencil is a convex average,
which also proves contraction in the maximum norm uniformly over grid sizes.
At the even-grid endpoint, the alternating mode has multiplier $-1$ and
remains bounded without decaying.
\item Write $a_k=-4s\sin^2(\pi k/N)\leq0$. The transformed equation gives
\[
 \wh u^{m+1}[k]=r_k\wh u^m[k],\qquad
 r_k=\frac{1+(1-\theta)a_k}{1-\theta a_k}.
\]
The denominator is positive, and $r_k\leq1$. The condition $r_k\geq-1$
is equivalent to $2+(1-2\theta)a_k\geq0$.
Thus the method is unconditionally stable for $\theta\geq1/2$.
For $\theta<1/2$, the exact fixed-grid condition is
$2s(1-2\theta)\mu_N\leq1$. The threshold is $1/2$, not $1/4$.
Program~\ref{listing-resolimplicite} evaluates any nonnegative integer
number of steps directly in Fourier coordinates.
\begin{listing}[H]
\PythonCode{listing-resolimplicite}
\caption{The one-dimensional theta method in Fourier coordinates.}
\label{listing-resolimplicite}
\end{listing}
\item On a square grid,
\[
 a_{kl}=-4s\bigl(\sin^2(\pi k/N)+\sin^2(\pi l/N)\bigr).
\]
For $\theta<1/2$, stability requires
$4s(1-2\theta)\mu_N\leq1$. In particular, the even-grid explicit bound
is $s\leq1/4$, not the one-dimensional bound $1/2$.
For $\theta\geq1/2$, stability is unconditional.
At $\theta=1/2$, $r_{kl}$ approaches $-1$ as $-a_{kl}$ grows:
high-frequency components alternate in sign while decaying very slowly.
At $\theta=1$, $r_{kl}=1/(1-a_{kl})$ approaches zero, so these components
are strongly damped. Program~\ref{listing-resolimplicite2d} implements both.
\begin{listing}[H]
\PythonCode{listing-resolimplicite2d}
\caption{The theta method on a periodic rectangular grid.}
\label{listing-resolimplicite2d}
\end{listing}
\end{enumerate}
\end{correction}

\begin{correction}{exo-unicite-eq-heat}
\begin{enumerate}
\item For $t>0$, absolute uniform convergence permits exchange of the integral
and the series:
\[
 p_t*f(x)=\sum_n e^{-2\pi^2n^2t}e^{2\pi\imath nx}
          \int_0^1f(y)e^{-2\pi\imath ny}\d y.
\]
On $t\geq\epsilon>0$, every differentiated series is dominated by a
polynomial in $|n|$ times $e^{-2\pi^2n^2\epsilon}$. Hence the convolution
is smooth and solves $u_t=\tfrac12u_{xx}$.
\item Two integrations by parts give $|c_n(f)|=O(|n|^{-2})$ for $n\ne0$.
Consequently,
\[
 \|u(t,\cdot)-f\|_\infty
 \leq\sum_n|c_n(f)|\,|e^{-2\pi^2n^2t}-1|\longrightarrow0
 \quad(t\downarrow0)
\]
by dominated convergence.
\item If $u(t_0,x_0)<0$, then $v=e^{\beta t}u$ attains a negative minimum
$\alpha$ on the compact cylinder $[0,t_0]\times S^1$.
Since the initial values are nonnegative, its minimizing time $t_1$ is
positive. At that point $v_{xx}\geq0$ and $v_t\leq0$; the latter also
holds at $t_1=t_0$ by the left derivative. But
$0\geq v_t=\beta\alpha+\tfrac12v_{xx}\geq\beta\alpha>0$
when $\beta<0$, a contradiction.
\item Apply positivity to smooth nonnegative initial functions, for which
continuity at $t=0$ was proved above. If $p_t$ were negative somewhere,
a nonnegative smooth bump supported near that point would give a negative
convolution value. Thus $p_t\geq0$. Its zeroth Fourier coefficient is one,
so its integral is one, and the integral triangle inequality gives
$\|p_t*f\|_\infty\leq\|f\|_\infty$.
For uniqueness, apply the minimum argument directly to $w$ and $-w$, where
$w$ is the difference of two real classical solutions with the same initial
values. Both have zero initial data, so $w=0$. For complex solutions apply
this to real and imaginary parts. This avoids assuming a representation
formula for an arbitrary solution.
\item Choose smooth periodic $f_j$ converging uniformly to $f$. Then
\[
 \|p_t*f-f\|_\infty
 \leq2\|f-f_j\|_\infty+\|p_t*f_j-f_j\|_\infty.
\]
First choose $j$ large, then let $t\downarrow0$. The first term is uniformly
small and the second tends to zero. This proves existence with the required
initial continuity, and the preceding maximum argument proves uniqueness.
\end{enumerate}
\end{correction}

\begin{correction}{exo-potential-electric-3d}
The stencil has coefficient $-6/h^2$ at the center and $1/h^2$ at each
of the six nearest neighbors. Each prescribed boundary neighbor is subtracted
from the right-hand side with factor $h^{-2}$. Odd extension in every
coordinate makes the boundary values zero and yields a $2N$-periodic array.
Its symbol is
\[
 -\frac4{h^2}\left(\sin^2\frac{\pi k}{2N}
                  +\sin^2\frac{\pi l}{2N}
                  +\sin^2\frac{\pi m}{2N}\right).
\]
Divide the transformed right-hand side by this symbol away from zero and
set the zero coefficient to zero. Invert the transform, retain the interior,
and restore the prescribed boundary. Program~\ref{listing-antisymmetrise}
constructs the odd extension in any dimension. The general solver in
Program~\ref{fish-listing} likewise works in two or three dimensions.
\begin{listing}[H]
\PythonCode{listing-antisymmetrise}
\caption{Odd extension along every coordinate axis.}
\label{listing-antisymmetrise}
\end{listing}
\end{correction}

\begin{correction}{exo-formulation-matrix-fish}
\begin{enumerate}
\item Multiplication by $T$ on the left takes second differences in the first
index; multiplication on the right takes them in the second. Their sum is
the five-point stencil with zero boundary values.
\item A boundary value is a known contribution to that stencil, so move it
to the right with coefficient $-h^{-2}$. Summing the contributions at corners
of the interior grid gives precisely $\wt F$ in the main text.
\item With $V_j[i]=\sin(\pi ij/N)$ and zero values at $i=0,N$,
\[
 (TV_j)[i]=\lambda_jV_j[i],\qquad
 \lambda_j=-4h^{-2}\sin^2(\pi j/(2N)).
\]
These eigenvalues are negative and distinct for $1\leq j<N$.
Let $U_0=V^{-1}\wt UV$ and $F_0=V^{-1}\wt FV$. Then
$(\lambda_i+\lambda_j)(U_0)_{ij}=(F_0)_{ij}$, so division determines
$U_0$ uniquely; no denominator vanishes. Finally $\wt U=VU_0V^{-1}$.
\item Symmetry of $T$ gives orthogonality between distinct eigenvectors,
and the identity $\sum_{i=1}^{N-1}\sin^2(\pi ij/N)=N/2$ gives
$V^TV=(N/2)I$. Since $V=V^T$, its inverse is $(2/N)V$.
Thus $V$ is \emph{scaled orthogonal}; $Q=\sqrt{2/N}V$ is orthogonal.
The product $Vx$ is the unnormalized sine transform $\Ss x$.
\item For the length-$2N$ odd extension
$x^o=(0,x[1],\ldots,x[N-1],0,-x[N-1],\ldots,-x[1])$,
\[
 \widehat{x^o}[k]
 =\sum_{j=1}^{N-1}x[j](e^{-\pi\imath jk/N}-e^{\pi\imath jk/N})
 =-2\imath\Ss(x)[k].
\]
Hence $\Ss x=\widehat{x^o}[1:N]/(-2\imath)$ and
$\Ss^{-1}=(2/N)\Ss$. The identity holds for complex data as well; taking
an imaginary part alone would discard information in that case.
Program~\ref{listing-transfo-sinus} preserves complex input.
\begin{listing}[H]
\PythonCode{listing-transfo-sinus}
\caption{The unnormalized sine transform, preserving complex data.}
\label{listing-transfo-sinus}
\end{listing}
For an interior matrix $F$, applying the transform in both coordinates gives
$\Ss_2(F)=VFV$, with inverse $(2/N)^2\Ss_2$.
The doubled odd array satisfies
$\widehat{F^o}[k,l]=-4(VFV)_{kl}$ on the interior frequencies.
Thus both constructions divide by the same eigenvalue sums and reconstruct
the same solution. Program~\ref{listing-transfo-sinus-2d} applies the
one-dimensional transform along the two axes.
\begin{listing}[H]
\PythonCode{listing-transfo-sinus-2d}
\caption{A two-dimensional sine transform.}
\label{listing-transfo-sinus-2d}
\end{listing}
\end{enumerate}
\end{correction}

\begin{correction}{exo-lissage-image}
Program~\ref{listing-calculfiltre} samples a Gaussian and normalizes its
finite sum to one. Its default radius is $\lceil4\sigma\rceil$ pixels.
Program~\ref{listing-appliquefiltre} reads a grayscale image and convolves
a reflected extension with this kernel, retaining the valid output.
Reflection preserves constant images and avoids an artificial dark border.
The noise is reduced at the cost of some edge and texture detail.
\begin{listing}[H]
\PythonCode{listing-calculfiltre}
\caption{A finite Gaussian kernel with unit sum.}
\label{listing-calculfiltre}
\end{listing}
\begin{listing}[H]
\PythonCode{listing-appliquefiltre}
\caption{Gaussian image smoothing with reflected boundaries.}
\label{listing-appliquefiltre}
\end{listing}
\end{correction}

\begin{correction}{exo-correlation-2d}
\begin{enumerate}
\item Expanding the square gives
\[
 d(f,g)[u,v]^2=E_f[u,v]+\|g\|_2^2-2\Corr(f,g)[u,v].
\]
If $E_f$ is constant, minimizing the squared discrepancy is exactly equivalent
to \emph{maximizing} correlation. If it varies little, this is only an approximation.
Full-template positions are essential here: otherwise the norm of the
visible part of $g$ also depends on the position.
\item Let $\check g[i,j]=g[-i,-j]$, with zero extension. Then
$\Corr(f,g)[u,v]=(f\star\check g)[u,v]$.
When the reflected template is stored as the ordinary array
$g[P-1-i,P-1-j]$, the corresponding full-convolution indices are
$(u+P-1,v+P-1)$. Padding to at least $N+P-1$ in both coordinates avoids
wraparound; extracting the valid part gives all complete-template positions.
\item Center the patch and template, writing their vectorized forms as
$a$ and $b$. For nonzero $a,b$,
$\rho=\langle a,b\rangle/(\|a\|_2\|b\|_2)$.
Cauchy--Schwarz gives $|\rho|\leq1$; equality $\rho=1$ means
$a=cb$ for some $c>0$, not necessarily $a=b$.
Offsets and positive rescalings of either intensity leave $\rho$ unchanged;
a negative rescaling of one image reverses its sign. A constant template
or patch has zero variance, so normalized correlation is undefined there.
\item Since $g_0=g-\mu_g$ has zero sum, the numerator is
$\sum_{i,j}f[u+i,v+j]g_0[i,j]$. This is one FFT-based correlation.
The patch variance is $s_2-s_1^2/P^2$, and the template variance
$\sum_{i,j}g_0[i,j]^2$ is computed once.
The stated recurrence follows by inclusion--exclusion of four overlapping
windows; the four corner terms must all contain the power $k$.
A summed-area table gives the equivalent formula
\[
 s_k[u,v]=I_k[u+P,v+P]-I_k[u,v+P]-I_k[u+P,v]+I_k[u,v],
\]
where $I_k[a,b]=\sum_{i<a,j<b}f[i,j]^k$.
Two cumulative sums construct $I_k$ in $O(N^2)$ work, after which each
window costs four lookups. Program~\ref{sliding-sum-listing} implements this.
\begin{listing}[H]
\PythonCode{sliding-sum-listing}
\caption{Sliding sums from a summed-area table.}
\label{sliding-sum-listing}
\end{listing}
Program~\ref{listing-correlation-normalized} combines these sums with FFT
correlation, for total cost $O(N^2\log N)$. It uses the Frobenius norm of
the template, subtracts a global image offset to reduce cancellation, and
marks zero-variance patches with \texttt{NaN}. Tiny negative computed
variances are clipped to zero before square roots; this corrects roundoff,
not the mathematical definition.
\begin{listing}[H]
\PythonCode{listing-correlation-normalized}
\caption{Normalized correlation at complete-template positions.}
\label{listing-correlation-normalized}
\end{listing}
\end{enumerate}
\end{correction}

\begin{correction}{exo-fft-rotation}
\begin{enumerate}
\item For signed frequencies $k,l$,
\[
 \widehat{T_vf}[k,l]=e^{-2\pi\imath(kv_1+lv_2)/N}\wh f[k,l].
\]
For integer shifts, this follows by changing indices. For fractional shifts,
it defines translation of the chosen trigonometric interpolant. Using
unsigned frequencies instead would select a different interpolation between
samples. Odd sizes avoid the unmatched real Nyquist mode of even sizes.
Program~\ref{listing-ffttranslation} implements the formula.
\begin{listing}[H]
\PythonCode{listing-ffttranslation}
\caption{Fractional translation using signed Fourier frequencies.}
\label{listing-ffttranslation}
\end{listing}
For an $x$-shear, translate each fixed-$y$ column by $\lambda y$; for a
$y$-shear, translate each fixed-$x$ row by $\lambda x$. Centered coordinates
place the center of the image at the origin.
Program~\ref{listing-ffttransvec} implements either axis.
\begin{listing}[H]
\PythonCode{listing-ffttransvec}
\caption{A centered shear along either image coordinate.}
\label{listing-ffttransvec}
\end{listing}
\item Multiply the three matrices with
$a=-\tan(\theta/2)$ and $b=\sin\theta$. The diagonal entries are
$1+ab=\cos\theta$, the lower-left entry is $b$, and the upper-right entry
is $2a+a^2b=-\sin\theta$. Since pullbacks compose in the corresponding
matrix order,
$R_\theta=S_a^{(x)}S_b^{(y)}S_a^{(x)}$ for continuous functions.
\item Each discrete line translation is unitarily diagonalized by the DFT,
with unit-modulus multipliers. Its inverse negates the shift, so each shear
and their product are invertible. Reversing the three steps and negating
their parameters gives the inverse; the symmetric three-step pattern makes
it the algorithm with angle $-\theta$.
\begin{listing}[H]
\PythonCode{listing-fftrotation}
\caption{Rotation by three discrete shears.}
\label{listing-fftrotation}
\end{listing}
Nevertheless, intermediate sheared functions do not generally belong to
the original finite tensor-product trigonometric space. Resampling and
periodic wrapping therefore prevent these discrete matrices from forming
an exact representation of all planar rotations:
$R_\alpha R_\beta\ne R_{\alpha+\beta}$ in general.
Taking real parts on even grids or cropping intermediate results can also
destroy invertibility. The odd-grid implementation preserves real data
without discarding a Nyquist component, and preserves complex data as complex.

Padding moves wraparound away from the image of interest; it does not make
band-limited interpolation exact for arbitrary continuous images.
Program~\ref{listing-applies-rotation} supplies a white margin for a small
rotation gallery. Larger angles should be split into bounded steps or
reduced using quarter-turn permutations to avoid very large shears.
The arithmetic cost is $O(N^2\log N)$ for a square padded grid.
\begin{listing}[H]
\PythonCode{listing-applies-rotation}
\caption{A rotation gallery with a padded white margin.}
\label{listing-applies-rotation}
\end{listing}
\end{enumerate}
\end{correction}

\begin{correction}{exo-creation-low-pass-filter}
\begin{enumerate}
\item Form the real even sampled response $H[k]=\mathbf1_{|\omega_k|<\pi/2}$,
with half-values at any cutoff samples, and take its inverse DFT.
The resulting circular kernel is real and symmetric, and its sum is $H[0]=1$.
\begin{listing}[H]
\PythonCode{listing-low-pass-filter}
\caption{A sharp sampled low-pass response and its inverse transform.}
\label{listing-low-pass-filter}
\end{listing}
\item Assign kernel indices to $-(N-1)/2,\ldots,(N-1)/2$, then embed them
modulo a larger plotting length. Padding the negative-index terms at the
wrong end would change the response. The resulting FFT samples
$G(\omega)=\sum_jg[j]e^{-\imath j\omega}$, which agrees with $H$ at the
original frequencies but oscillates between them near the cutoff.
Program~\ref{listing-drawing-filter-zero-padding} returns the frequencies
and the complex response; for this symmetric real kernel it is real up to roundoff.
\begin{listing}[H]
\PythonCode{listing-drawing-filter-zero-padding}
\caption{Dense frequency-response evaluation with centered kernel indices.}
\label{listing-drawing-filter-zero-padding}
\end{listing}
\item Let $r=|\omega|$, $a=(1-\epsilon)\pi/2$, and
$b=(1+\epsilon)\pi/2$. Use response one for $r\leq a$, zero for $r\geq b$,
and $\tfrac12[1+\cos(\pi(r-a)/(b-a))]$ in the transition.
At $\epsilon=0$ use the sharp-cutoff definition separately.
Increasing the transition width reduces the long oscillatory tails of the
impulse response. The endpoint $\epsilon=1$ is included without creating
a negative array length.
\begin{listing}[H]
\PythonCode{configurable-filter-listing}
\caption{A raised-cosine transition, including both endpoint parameters.}
\label{configurable-filter-listing}
\end{listing}
\end{enumerate}
\end{correction}

\begin{correction}{exo-iterations-integers}
\begin{enumerate}
\item The rounded update is
$a_j^{m+1}=2\lceil(a_j^m+a_{j-1}^m)/4\rceil$.
Its minimum is nondecreasing and its maximum nonincreasing, since both are
even. If the current minimum is $q$, an updated count equals $q$ exactly
when both contributing counts equal $q$. Thus its new minimum set is the
intersection of the old minimum set and its one-step translate.
A nonempty proper subset of a cycle cannot be invariant under a one-step
translation, so its cardinality strictly decreases while the minimum stays
fixed. Within at most $n$ steps, either consensus is reached or the minimum
increases by at least two. It is bounded above by the initial maximum,
so consensus occurs after finitely many steps. This argument excludes the
already-constant case before claiming strict decrease.
\item Without rounding, the kernel is $(1/2,1/2,0,\ldots,0)$.
Its nonconstant Fourier multipliers have modulus less than one, so the limit
is the constant vector with value $n^{-1}\sum_ja_j^0$.
\item With half sent to each neighbor, the multiplier is
$\cos(2\pi k/n)$. For odd $n$, its modulus is less than one except at zero,
so the fractional process converges to the mean. For even $n$, the
alternating mode has multiplier $-1$: alternating counts $0,2,0,2,\ldots$
are exchanged at every step and do not converge, even with rounding.
For the rounded process on an odd cycle, a site remains at the current
minimum exactly when both its neighbors were at that minimum.
If the minimum set failed to shrink, it would be invariant under translation
by two. Since two generates an odd cycle, a proper nonempty minimum set
cannot have this property. The preceding finite-consensus argument therefore
applies to odd $n$ as well.
\end{enumerate}
\end{correction}

\begin{correction}{exo-algo-karatsuba}
\begin{enumerate}
\item The terms are $R_0=P_0Q_0$, $R_2=P_1Q_1$, and
$R_1=P_0Q_1+P_1Q_0$.
\item Compute
$R_1=(P_0+P_1)(Q_0+Q_1)-R_0-R_2$.
The cross term needs one new product, two additions, and two subtractions.
\item After padding to length $n=2^p$, the recurrence is
$C(n)=3C(n/2)+O(n)$, hence
\[
 C(n)=O\!\left(\sum_{j=0}^{p-1}3^j\frac n{2^j}+3^p\right)
     =O(n^{\log_2 3}).
\]
Padding arbitrary lengths increases them by less than a factor two and
preserves the bound. Programs~\ref{listing-karatsuba} and \ref{listing-add}
use Python lists, so integer coefficients retain arbitrary precision.
The algebra also works over other coefficient rings supporting addition,
subtraction, and multiplication. There is no universal crossover degree:
benchmark comparable implementations and include conversion and precision costs.
\begin{listing}[H]
\PythonCode{listing-karatsuba}
\caption{Karatsuba multiplication with exact Python integer coefficients.}
\label{listing-karatsuba}
\end{listing}
\begin{listing}[H]
\PythonCode{listing-add}
\caption{Addition of coefficient vectors of different lengths.}
\label{listing-add}
\end{listing}
\end{enumerate}
\end{correction}

\begin{correction}{exo-spline-filtering}
\begin{enumerate}
\item The interpolation equations are $u_d=a*\psi_d$.
Since $\psi_d$ is finitely supported, its transform $D=\widehat{\psi_d}$
is continuous and periodic. Nonvanishing therefore implies
$\min_{[-\pi,\pi]}|D|>0$, so multiplication by $D$ is a boundedly
invertible operator on $L^2$ of the circle. This proves the unique solution
$\widehat a=\widehat{u_d}/D$.
Let $b$ have Fourier series $1/D$ and set
$\varphi(x)=\sum_kb[k]\psi(x-k)$. At each $x$, this sum has finitely many
nonzero terms; moreover $b*\psi_d=\delta_0$, so $\varphi(j)=\delta_{j0}$.
Then $a=b*u_d$ gives the direct form with kernel $\varphi$.
The inverse filter generally has infinitely many coefficients, so compact
support of $\psi$ need not imply compact support of $\varphi$.
\item Repeated convolution gives
$\Supp(\beta^n)=[-(n+1)/2,(n+1)/2]$ (as a closed support).
Direct interpolation requires $\beta^n(j)=\delta_{j0}$.
This holds for $n=0$ and $n=1$, giving piecewise constant and piecewise linear
interpolation. It fails for $n\geq2$; these degrees require a prefilter.
\item Put $r=n+1$. The continuous convolution theorem gives
\[
 B_r(\xi):=\widehat{\beta^n}(\xi)
          =\left(\frac{\sin(\xi/2)}{\xi/2}\right)^r.
\]
For $n\geq1$, Poisson summation, with absolutely convergent frequency tails,
gives
\[
 D_n(\xi)=\sum_{k\in\ZZ}\beta^n(k)e^{-\imath k\xi}
         =\sum_{m\in\ZZ}B_r(\xi+2\pi m).
\]
There is no factor $2\pi$ in this identity. By evenness and periodicity it
suffices to take $0\leq\xi\leq\pi$.
If $r$ is even, every summand is nonnegative and $B_r(\xi)>0$.
If $r\geq3$ is odd, put $x=\xi/(2\pi)$. For $0<x\leq1/2$, pairing the
terms $m=j$ and $m=-j$ gives
\[
 D_n(\xi)=B_r(\xi)+\left(\frac{\sin(\pi x)}\pi\right)^r
 \sum_{j=1}^\infty(-1)^{j+1}
       \left((j-x)^{-r}-(j+x)^{-r}\right).
\]
The parenthesized quantities are positive and decrease to zero: their
derivative as a function of $j$ is negative. The alternating sum is
nonnegative. Thus in both parity cases
$D_n(\xi)\geq B_r(\xi)>0$ on $[-\pi,\pi]$; at zero both sides equal one.
For $n=0$, the sampled sequence is simply $\delta_0$, so $D_0=1$.
\item The transform of the cardinal kernel is $H_r=B_r/D_n$.
For $|\eta|\leq\pi$ and $m\in\ZZ$,
\[
 |H_r(\eta+2\pi m)|
 \leq\left|\frac{\eta}{\eta+2\pi m}\right|^r,
\]
with the $m=0$ ratio interpreted as one. This follows from
$D_n(\eta)\geq B_r(\eta)>0$ and the ratio of the two sinc powers.
For $|\eta|<\pi$, every ratio with $m\ne0$ is strictly less than one, and
the shifted series shows $D_n(\eta)/B_r(\eta)\to1$ as $r\to\infty$.
Consequently $H_r(\eta)\to1$ and $H_r(\eta+2\pi m)\to0$ for $m\ne0$,
except at a set of endpoints of measure zero.

For $r\geq2$ and $m\ne0$, the squared terms are bounded by
$(2|m|-1)^{-4}$. This bound is summable in $m$ and integrable over
$\eta\in[-\pi,\pi]$, while $|H_r(\eta)|\leq1$ on the central interval.
Dominated convergence on these intervals therefore proves
$H_r\to\mathbf1_{[-\pi,\pi]}$ in $L^2(\RR)$.
Since $\widehat{\sinc}=\mathbf1_{[-\pi,\pi]}$ in our convention,
Plancherel gives $\beta^n_{\mathrm{card}}\to\sinc$ in $L^2$.
The argument keeps the alternating signs for even spline degrees; dropping
them would give a different cardinal kernel.

For the cubic spline, $D_3(\xi)=(2+\cos\xi)/3$ and its inverse filter is
$b[k]=\sqrt3(\sqrt3-2)^{|k|}$. Hence
\[
 \beta^3_{\mathrm{card}}(x)=\sqrt3\sum_{k\in\ZZ}
              (\sqrt3-2)^{|k|}\beta^3(x-k).
\]
It has unbounded, exponentially decaying tails, whereas sinc decays only
algebraically. The limiting cardinal interpolation uses sinc translates,
as in Shannon's sampling theorem.
\item For a finite vector, division by sampled values of $D_n$ followed by
an inverse FFT computes the \emph{periodic} prefilter exactly for that finite
model. It approximates the infinite-grid problem only after a boundary model
and sufficient extension have been chosen. The principal issue is the
unintended periodic boundary condition, rather than an unavoidable failure
of FFT convolution. Recursive causal and anticausal filters exploit the
roots of the spline symbol; Exercise~\ref{exo-calcul-spline-iir} develops
this method for cubic splines. Natural and not-a-knot finite-interval splines
solve different boundary-value problems and should not be conflated with
periodic or infinite-grid interpolation.
\end{enumerate}
\end{correction}
